\documentclass{article}
\usepackage[dblblindworkshop]{neurips_2026}
\workshoptitle{AI Foundations for Power Grids}

\usepackage[utf8]{inputenc}
\usepackage[T1]{fontenc}
\usepackage{hyperref}
\usepackage{url}
\usepackage{booktabs}
\usepackage{amsfonts}
\usepackage{amsmath}
\usepackage{nicefrac}
\usepackage{microtype}
\usepackage{xcolor}
\usepackage{graphicx}

\newcommand{\arr}{ARR}
\newcommand{\todo}[1]{\textcolor{red}{[TODO: #1]}}

\title{Adaptive Reliability Reserves: Learning Risk-Adjusted Backup for Residential Battery Fleets}

\begin{document}

\maketitle

\begin{abstract}
Residential battery virtual power plants promise each household outage backup while monetizing the same batteries in wholesale markets. Industry practice resolves this tension with a fixed state-of-charge floor (typically 20\%) plus a binary weather-alert ``storm mode.'' We ask whether the backup reserve should instead be a risk-adjusted sequential decision. We formalize the customer promise as a per-home reliability-risk budget---expected unserved backup energy per year---and compare fixed floors, alert rules, quantile reserves, risk-priced model-predictive control, and a learned policy that sets fleet-level dispatch and reserve targets above a feasibility-preserving disaggregation layer. Evaluation uses a public replica of a Texas home-battery fleet: NREL ResStock homes, ERCOT prices (2019--2026), and ORNL EAGLE-I county outage records (2014--2025). \todo{headline numbers: frontier improvement at equal realized reliability; behavior under Winter Storm Uri; fleet-size/composition shift.} Price spikes and mass outages are nearly disjoint in ERCOT (2021--2025: 107 spike days, 28 mass-outage days, 4 joint---all Winter Storm Uri), so hazard-conditioned reserves are close to free most of the year; we characterize where they fail.
\end{abstract}

\section{Introduction}
\todo{Tension: dual-use storage. Base-style fixed 20\% floor + storm hold; Tesla Storm Watch; Sunrun 20\%. Why always 20\%? Contributions: (1) reserve as learned allocation of a per-home reliability budget, nesting industry rules and quantile reserves; (2) HomeFleet-ERCOT, a public benchmark coupling real prices, real outage records, and realistic homes; (3) evaluation under market, outage, and fleet shift with failure analysis.}

\section{Problem: backup promises as a risk budget}
\todo{Firm margin (A\&Z); battery dynamics; outage process from EAGLE-I via quantile matching; unserved energy; per-home budget CMDP; multiplier = implied VoLL; special cases: fixed floor, storm mode, quantile tier, chance-constrained MPC.}

\section{Controllers}
\todo{DP/MPC family with shared forecasts; risk-priced MPC (same objective as RL); PPO + disaggregation layer; action = aggregate power + reserve multiplier; feasibility by construction.}

\section{Benchmark and protocol}
\todo{Data sources and splits (train $\le$2023, val 2024, test 2025--26, event holdouts); the 2$\times$2 spike-outage event design; metrics: margin/home-week, capture vs oracle, CVaR$_5$, ride-through and 5-hour promise rates, unserved kWh; common random numbers.}

\section{Results}
\todo{Fig 2 orthogonality; Fig 3 margin-reliability frontier; Fig 4 event panels (Uri, Beryl, Aug-2023 heat); Fig 5 shift (years, RTC+B, fleet size/composition, dropout); Table tails + failure taxonomy. Early findings to confirm: risk-priced MPC beats fixed floor in every spike window at equal reliability; storm-watch loses money in heat waves but is the best hedge in Uri; forecast-based controllers all lose money in Uri (REP exposure).}

\section{Limitations and deployment notes}
\todo{County-level hazard; price-taker; AMY2018 loads not cold-amplified; ADER caps; sim-to-real (6.7\% precedent); what an operator would need before eyes-off.}

\begin{ack}
\todo{remove for submission (double-blind).}
\end{ack}

\section*{LLM usage disclosure}
\todo{Statement per NeurIPS 2026 policy: LLM-assisted coding and editing; all citations and numbers verified by the authors.}

\bibliographystyle{plainnat}
\bibliography{refs}

\appendix
\section{Domain checklist}
\begin{enumerate}
\item \textbf{Physics feasibility.} \todo{SOC/inverter/backup feasibility by construction (disaggregation layer); AC power-flow verification on an NREL SMART-DS Austin feeder; violation statistics.}
\item \textbf{Out-of-distribution evaluation.} \todo{held-out years incl.\ post-RTC+B prices; event holdouts excluded from training; unseen homes; fleet size/composition shift.}
\item \textbf{Failure modes.} \todo{Uri panel; tail stats (CVaR$_5$, worst week); failure taxonomy by event class.}
\item \textbf{System-level effect.} \todo{hazard-model quality $\to$ downstream margin/reliability (value-of-information bracket: alerts-only vs oracle hazard).}
\item \textbf{Data and code.} \todo{all-public data; anonymized repo link; full release on acceptance.}
\end{enumerate}

\section{Extended results}
\todo{per-window tables; degradation-model sensitivity; ablations.}

\end{document}
